\section{Preliminaries and Setup}
\label{sec:preliminaries and setup}
\subsection{Preliminaries}
\paragraph{On-policy distillation.}
Given a question $x$ from the training set $\mathcal D$, OPD samples a response $y=(y_1,\ldots,y_n)$ from the current student $p_{\theta_t}$. A fixed teacher $q$ supplies a next-token distribution at each student-generated prefix $c_i=(x,y_{<i})$. The student is updated to reduce its discrepancy from the teacher on these prefixes, and subsequent rollouts use the updated student.

We use reverse KL on a vocabulary support $S(c)$. Full-vocabulary reverse KL takes $S(c)=\mathcal V$, where $\mathcal V$ is the vocabulary; our experiments use the student's Top-16 support. Define $\tilde p_\theta(v\mid c)=p_\theta(v\mid c)/\sum_{u\in S(c)}p_\theta(u\mid c)$ and $\tilde q$ analogously. The position loss is
\begin{equation}
\ell_\theta(c)=\sum_{v\in S(c)}\tilde p_\theta(v\mid c)
\log\frac{\tilde p_\theta(v\mid c)}{\tilde q(v\mid c)}.
\label{eq:loss}
\end{equation}
For a sampled batch $\mathcal B$ with response lengths $n_b$, we minimize
\begin{equation}
L_{\mathcal B}(\theta)=\frac1{N_{\mathcal B}}\sum_{b\in\mathcal B}\sum_{i=1}^{n_b}\ell_\theta(c_{bi}),
\qquad N_{\mathcal B}=\sum_{b\in\mathcal B}n_b.
\label{eq:training_loss}
\end{equation}
Responses, student-selected support, and teacher targets are fixed during each update.

Further more, to keep notation consistent throughout the paper, we define the following quantities. At checkpoint $t$, write $\ell_t=\ell_{\theta_t}$ and $h_t(c)=h_{\theta_t}(c)\in\mathbb R^d$ for the loss and final hidden state. The output head $W_t$ produces logits $z_t(c)=W_t h_t(c)$. We denote the individual-position hidden gradient by $g_{h,t}(c)=\nabla_h\ell_t(c)$, the training parameter gradient by $g_T=\nabla_\theta L_{\mathcal B}$, and the realized optimizer step by $\Delta\theta=\theta^+-\theta$, where $\theta^+$ denotes the parameters after the update. Checkpoint subscripts are omitted when clear from context. We study OPD training dynamics from four perspectives: the origins and changes of the learning signal, adaptation in representation space, transfer to unseen responses, and training efficiency.

\subsection{Setup}

\paragraph{Experimental setup.}
We distill Qwen3-1.7B-Base and Qwen3-4B-Base from their respective
post-trained teachers, Qwen3-1.7B and Qwen3-4B. After filtering,
DAPO14k contains 13,534 math questions. For the paired experiments,
we reserve 128 questions for validation and train on the remaining
13,406. We use AdamW with a constant learning rate of $10^{-5}$ for
full-parameter training (FULL). We evaluate on MATH500, HMMT25,
HMMT26, AIME26, and Minerva. For each benchmark, we report Mean@8,
the percentage of correct answers averaged over eight sampled responses
per question and then over questions. We compute the overall score
as the unweighted mean of the five benchmark scores.
Appendix~\ref{app:settings} gives the full experimental settings.